\documentclass{article}
\usepackage{graphicx} % Required for inserting images
\usepackage[utf8]{inputenc}
\usepackage[T1]{fontenc} % acentos em verbatim e < > corretos
\usepackage{mathptmx}     % Times New Roman (texto e fórmulas)
\usepackage{courier}      % Courier para código (vetorial, combina com a Times)
\setlength{\emergencystretch}{2em} % evita linhas a passar a margem com nomes longos
\usepackage[portuguese]{babel}
\usepackage{booktabs}
\usepackage{array}
\usepackage{multirow}
\usepackage{geometry}
\geometry{a4paper, margin=2.5cm}
\usepackage{hyperref}

\title{Bolsa Investigação - Centro Algoritmi\\ Avaliação de LLMs com RAG sobre as atas da Câmara Municipal de Barcelos}
\author{
    Beatriz Miranda \\
    Tiago Martins
}
\date{2026}

\begin{document}

\maketitle
\newpage

\tableofcontents
\newpage

\section{Introdução}
\label{sec:intro}

Este relatório documenta o projeto de avaliação comparativa de modelos de linguagem (LLMs) para extração de informação das atas das reuniões da Câmara Municipal de Barcelos, utilizando um pipeline RAG (Retrieval-Augmented Generation) padronizado. O objetivo principal é comparar o desempenho de diferentes modelos na tarefa de responder a perguntas sobre as deliberações municipais.

O estudo tem dois eixos de experiência independentes. No primeiro, fixa-se o pipeline RAG e varia-se apenas o modelo de geração, para comparar modelos. No segundo, fixa-se o modelo e varia-se uma decisão do pipeline de cada vez (OCR, chunking, pesquisa ou prompt), para medir o efeito de cada decisão.

\subsection{Contexto e Motivação}

As atas das reuniões de câmara contêm informação estruturada mas complexa, com:
\begin{itemize}
    \item Textos em português europeu com terminologia jurídica e administrativa
    \item Datas escritas por extenso (ex: ``aos trinta dias do mês de janeiro de dois mil e vinte e cinco'')
    \item Estrutura hierárquica: atas $\rightarrow$ propostas $\rightarrow$ deliberações
    \item Informação distribuída por múltiplos documentos
\end{itemize}

A avaliação de LLMs neste contexto permite entender:
\begin{itemize}
    \item Qual a capacidade de cada modelo lidar com português europeu administrativo
    \item Como o tamanho do modelo afeta o desempenho em tarefas de extração
    \item A eficácia de modelos genéricos vs. especializados em português
    \item A resistência à alucinação face a perguntas sem resposta
\end{itemize}

\section{Corpus de Atas}
\label{sec:corpus}

O corpus é composto por 23 atas das reuniões da Câmara Municipal de Barcelos de 2025, obtidas em PDF no sítio do Município:
\begin{itemize}
    \item atas n.º 1 a 20, de janeiro a setembro de 2025 (a ata n.º 18 não existe), incluindo uma reunião extraordinária (ata n.º 9, de 22 de abril);
    \item atas n.º 1 a 4 do mandato 2025--2029, de novembro e dezembro de 2025.
\end{itemize}

Duas características do corpus condicionaram o desenho do pipeline:
\begin{itemize}
    \item \textbf{Numeração repetida.} Como a numeração recomeça com o novo mandato, existem duas atas com cada um dos números 1 a 4. Por isso, as atas são identificadas pela data da reunião (AAAA-MM-DD), tanto nas perguntas (campo \texttt{ata\_fonte}) como na métrica de recuperação.
    \item \textbf{Ata digitalizada.} A ata n.º 5 (3 de março de 2025) é quase toda uma imagem digitalizada: só uma página tem texto extraível. É um dos motivos para o pipeline aplicar OCR antes da extração de texto (Secção~\ref{subsec:ocr}).
\end{itemize}

\subsection{Extração de Metadados}

O número, a data e o tipo de cada reunião são extraídos do título da ata (\textit{``Ata n.º X da Reunião Ordinária \ldots{} realizada em \ldots''}). A extração exigiu tratar três particularidades:
\begin{itemize}
    \item o número da ata aparece numa linha isolada no PDF e seria removido pela limpeza de números de página, pelo que os metadados são lidos do texto antes da limpeza;
    \item a posição do título no texto extraído depende do extrator (o PyMuPDF coloca-o depois do primeiro parágrafo), pelo que o título é procurado e não assumido no início;
    \item o tipo de reunião é lido apenas do título, já que o corpo das atas ordinárias menciona frequentemente reuniões extraordinárias (por exemplo, na aprovação de atas anteriores).
\end{itemize}
Quando o título não é legível, os metadados são obtidos do nome do ficheiro.

\section{Arquitetura do Pipeline RAG}
\label{sec:pipeline}

\subsection{Visão Geral}

O pipeline segue uma arquitetura modular em quatro etapas:

\begin{center}
\begin{tabular}{c}
\toprule
PDFs das atas $\rightarrow$ OCR $\rightarrow$ ingest.py $\rightarrow$ chunks.jsonl \\
$\rightarrow$ index.py $\rightarrow$ ChromaDB (embeddinggemma) \\
$\downarrow$ \textit{top-k excertos, iguais para todos os modelos} \\
questions.jsonl $\rightarrow$ run.py $\rightarrow$ answers\_<pipeline>\_\_<modelo>.jsonl \\
$\downarrow$ \\
evaluate.py $\rightarrow$ CSVs (resultados) \\ \bottomrule
\end{tabular}
\end{center}

Cada pipeline é uma configuração RAG completa, definida no ficheiro \texttt{config.yaml} com as secções OCR, chunking, embeddings, pesquisa, geração e prompt. Os modelos ficam fora desta estrutura, o que permite fixar um deles e variar apenas o pipeline.

\subsection{Decisões de Desenho}
\label{subsec:decisoes}

As decisões de desenho foram tomadas com o objetivo de isolar a variável em estudo e garantir uma comparação justa. As principais decisões do pipeline de referência (\texttt{baseline}) são:

\begin{table}[h!]
\centering
\caption{Decisões de desenho do pipeline RAG}
\label{tab:decisoes}
\begin{tabular}{>{\raggedright\arraybackslash}p{7cm}>{\raggedright\arraybackslash}p{8cm}}
\toprule
\textbf{Decisão} & \textbf{Justificação} \\ \midrule
OCR de todas as páginas (OCRmyPDF/Tesseract, português, com correção de rotação) & Todas as atas chegam à extração pelo mesmo processo, incluindo as digitalizadas \\ \midrule
Recuperação fixa partilhada por todos os modelos & Isola a variável em estudo: apenas o gerador muda, permitindo comparação justa entre modelos \\ \midrule
Chunking por ``PROPOSTA N.º X'' + janela de 1500 caracteres com sobreposição de 200 & Cada deliberação fica num chunk coerente; evita misturar propostas diferentes no mesmo chunk \\ \midrule
Cabeçalho com metadados em cada chunk: [Ata n.º X | tipo | data] & As atas usam datas por extenso; o cabeçalho torna perguntas por data/n.º de ata recuperáveis \\ \midrule
Atas identificadas pela data da reunião, e não pelo número & A numeração recomeça em cada mandato (Secção~\ref{sec:corpus}) \\ \midrule
Embeddings multilingues: embeddinggemma (servidor Ollama da universidade), com prefixos distintos para perguntas e documentos & Modelo multilingue treinado para pesquisa, disponível na infraestrutura da universidade; o mesmo modelo é usado para indexar e pesquisar \\ \midrule
Top-k = 5 para todos os modelos & Recuperação idêntica para todos, eliminando variabilidade \\ \midrule
temperature = 0 & Respostas determinísticas e reprodutíveis \\ \midrule
Prompt obriga a responder ``A informação não consta das atas fornecidas.'' quando não há informação & Permite medir alucinação de forma objetiva e determinística \\ \midrule
Orquestrador/juiz RAGAS separado dos modelos avaliados (Gemma 4, 26B) & Evita que o juiz favoreça o próprio modelo, com capacidade superior e custo razoável \\ \midrule
Limpeza de texto: remoção de tracejados, números de página, hifenização & Normalização do texto extraído dos PDFs \\ \midrule
Extração de datas por extenso para formato ISO & Conversão automática de datas como ``trinta e um de março de dois mil e vinte e cinco'' para ``2025-03-31'' \\ \bottomrule
\end{tabular}
\end{table}

\subsection{OCR}
\label{subsec:ocr}

Todas as atas passam por OCR (OCRmyPDF com Tesseract, língua portuguesa) antes da extração de texto, incluindo as páginas que já têm texto, que são rasterizadas e reconhecidas de novo. Assim, todas as atas chegam ao modelo de embeddings e aos LLMs pelo mesmo processo. Os PDFs reconhecidos são guardados em cache, pelo que o OCR só é feito uma vez por configuração.

O OCR pode introduzir erros de reconhecimento, sobretudo em algarismos (valores, datas e números de proposta). O efeito do OCR na qualidade da recuperação e das respostas é, por isso, uma das variáveis a analisar.

\subsection{Pesquisa}
\label{subsec:pesquisa}

O pipeline de referência usa pesquisa semântica: a pergunta é convertida num embedding e o ChromaDB devolve os 5 excertos mais próximos pela semelhança do cosseno. As variantes de pesquisa (Secção~\ref{sec:pipelines}) acrescentam:
\begin{itemize}
    \item \textbf{Pesquisa híbrida.} A pesquisa semântica capta sinónimos e paráfrases, mas falha em nomes próprios, códigos e termos raros, onde a pesquisa por palavras-chave (BM25) funciona melhor. O BM25 é calculado em português, com redução das palavras à raiz (\textit{stemming} Snowball: ``aprovada'', ``aprovação'' e ``aprovar'' $\rightarrow$ ``aprov'') e sem palavras comuns. As duas listas são fundidas por \textit{Reciprocal Rank Fusion}, $\mathrm{score}(d) = \sum 1/(60 + \mathrm{posição}(d))$, que combina as posições e não os scores, já que os scores do cosseno e do BM25 não são comparáveis. Uma variante alternativa usa a soma ponderada dos scores normalizados (0,7 semântico e 0,3 BM25).
    \item \textbf{Filtro por data.} A pesquisa vetorial não é adequada a valores exatos como datas. Quando a pergunta indica a data de uma reunião existente no corpus, a pesquisa é restringida aos excertos dessa ata, através de um filtro sobre os metadados. Datas sem reunião correspondente (por exemplo, em perguntas com premissas falsas) são ignoradas.
    \item \textbf{Reranking.} Recuperam-se os 20 excertos mais próximos e um \textit{cross-encoder} multilingue (\texttt{mmarco-\allowbreak mMiniLMv2}), que lê a pergunta e o excerto em conjunto, reordena-os, ficando os 5 melhores.
\end{itemize}

\section{Seleção dos Modelos}
\label{sec:modelos}

\subsection{Critérios de Seleção}

Os modelos foram selecionados para comparar o desempenho de diferentes famílias e estratégias de afinação numa gama de tamanho semelhante. Esta abordagem permite:

\begin{itemize}
    \item Avaliar se diferentes famílias de modelos, com tamanhos comparáveis, apresentam desempenhos distintos no domínio das atas municipais
    \item Identificar quais os modelos que melhor lidam com o português europeu administrativo
    \item Comparar modelos abertos disponíveis no servidor da universidade
    \item Isolar o efeito da estratégia de afinação dentro da mesma família e escala
\end{itemize}

\subsection{Modelos a Avaliar}

Foram selecionados seis modelos entre 8 e 14 mil milhões de parâmetros, todos disponíveis na infraestrutura da universidade: Qwen~3~8B, Gervásio~8B PT-PT, AMALIA 0626 SFT~9B, AMALIA 0626 DPO~9B, Phi-4~14B e Ministral~3~14B. Esta é a gama compacta disponível mais próxima de 10B; o único par da mesma família é deliberado e permite uma comparação controlada entre estratégias de afinação.

Qwen funciona como referência generalista e Gervásio testa especialização em português europeu. AMALIA SFT e DPO permitem avaliar o efeito da estratégia de afinação mantendo a família e a escala constantes. Phi-4 e Ministral são comparados entre si na mesma escala de 14B. Os resultados serão sempre apresentados com o número de parâmetros: não se atribuirá uma diferença entre 8B e 14B apenas à família do modelo. Os identificadores exatos, endpoint e método de autenticação estão registados em \texttt{config.yaml}; cada membro usa a sua chave local \texttt{OLLAMA\_TOKEN}, que não é versionada.

\subsection{Orquestrador e Modelo Juiz}

O modelo \texttt{gemma4:26b} foi escolhido como orquestrador/juiz das métricas RAGAS. Tem 26B parâmetros, acima dos geradores de 8--14B, e não pertence à amostra avaliada. Desta forma, reduz-se o risco de auto-favorecimento sem recorrer ao custo de um modelo de 119B.

Neste projeto, ``orquestrador'' significa o avaliador que coordena o julgamento das respostas nas métricas RAGAS; não responde às perguntas do conjunto experimental. As respostas medidas são sempre produzidas apenas pelos seis modelos selecionados.

\section{Variantes do Pipeline}
\label{sec:pipelines}

Cada variante altera uma única decisão do pipeline de referência, para que as diferenças observadas sejam atribuíveis a essa decisão:

\begin{table}[h!]
\centering
\caption{Pipelines definidos em \texttt{config.yaml}}
\label{tab:pipelines}
\begin{tabular}{>{\raggedright\arraybackslash}p{3.5cm}>{\raggedright\arraybackslash}p{11cm}}
\toprule
\textbf{Pipeline} & \textbf{Alteração em relação ao \texttt{baseline}} \\ \midrule
baseline & Referência: OCR com correção de rotação, chunks por proposta, embeddinggemma, top-5 e prompt com citação de fontes \\ \midrule
ocr\_deskew & OCR com correção da inclinação das páginas \\ \midrule
chunks\_2500 & Chunks de 2500 caracteres com sobreposição de 300 \\ \midrule
top\_k\_10 & 10 excertos em vez de 5 \\ \midrule
hybrid\_baseline & Pesquisa híbrida (semântica + BM25 em português), fundida por RRF \\ \midrule
hybrid\_weighted & Pesquisa híbrida com soma ponderada dos scores (0,7/0,3) \\ \midrule
rerank & Reranking dos 20 primeiros excertos por um \textit{cross-encoder} multilingue \\ \midrule
filtro\_data & Filtro pela data da reunião, quando a pergunta a indica \\ \midrule
prompt\_sem\_fontes & Prompt sem obrigação de citar a ata no fim da resposta \\ \bottomrule
\end{tabular}
\end{table}

Os pipelines que só alteram a pesquisa ou o prompt produzem o mesmo OCR, os mesmos chunks e o mesmo índice que o \texttt{baseline}, pelo que reutilizam esses artefactos. Apenas \texttt{ocr\_deskew} e \texttt{chunks\_2500} exigem OCR ou indexação próprios.

\section{Métricas de Avaliação}
\label{sec:metricas}

\subsection{Métricas RAGAS}

O framework RAGAS (Retrieval-Augmented Generation Assessment) fornece métricas baseadas em LLM-juiz para avaliar tanto a recuperação como a geração:

\begin{table}[h!]
\centering
\caption{Métricas RAGAS implementadas}
\label{tab:ragas}
\begin{tabular}{>{\raggedright\arraybackslash}p{6cm}>{\raggedright\arraybackslash}p{5cm}>{\raggedright\arraybackslash}p{3cm}}
\toprule
\textbf{Métrica} & \textbf{O que mede} & \textbf{Componente} \\ \midrule
faithfulness & A resposta é suportada pelos excertos? (anti-alucinação) & Geração \\ \midrule
answer\_relevancy & A resposta responde mesmo à pergunta? & Geração \\ \midrule
factual\_correctness & Afirmações da resposta vs. referência (F1) & Geração \\ \midrule
llm\_context\_precision\_\allowbreak with\_reference & Os excertos úteis aparecem no topo? & Recuperação \\ \midrule
context\_recall & Os excertos cobrem a informação da referência? & Recuperação \\ \bottomrule
\end{tabular}
\end{table}

A métrica answer\_relevancy compara embeddings, calculados com o mesmo modelo usado na pesquisa (embeddinggemma).

\subsection{Métricas Próprias (Determinísticas)}

Para complementar o RAGAS e validar os seus resultados, foram implementadas métricas determinísticas que não dependem de um juiz LLM:

\begin{table}[h!]
\centering
\caption{Métricas próprias implementadas}
\label{tab:metricas-proprias}
\begin{tabular}{>{\raggedright\arraybackslash}p{3.5cm}>{\raggedright\arraybackslash}p{6.5cm}>{\raggedright\arraybackslash}p{4.5cm}}
\toprule
\textbf{Métrica} & \textbf{Definição} & \textbf{Avalia} \\ \midrule
valores\_corretos & Percentagem de valores\_esperados presentes na resposta & Precisão factual \\ \midrule
abstencao\_correta & Respondeu quando havia resposta / recusou quando não havia & Comportamento adequado \\ \midrule
alucinacao & Respondeu a uma pergunta sem resposta nas atas (deveria recusar) & Medida direta de alucinação \\ \midrule
falsa\_recusa & Disse ``não consta'' quando a resposta estava disponível & Sub-utilização do contexto \\ \midrule
retrieval\_hit@k & Alguma das atas corretas (identificadas pela data) está nos top-k excertos & Qualidade da recuperação \\ \midrule
latencia\_s & Tempo de resposta em segundos & Eficiência \\ \midrule
tokens\_out & Número de tokens gerados & Verbosidade/custo \\ \bottomrule
\end{tabular}
\end{table}

Uma pergunta é considerada sem resposta quando a sua resposta de referência é uma recusa (``não consta'', ``não é possível'') ou um pedido de clarificação, independentemente do seu tipo. Este critério é necessário porque vários tipos de robustez (por exemplo, \texttt{resistencia\_a\_entidade\_inexistente}) têm como resposta correta a recusa. Na resposta do modelo, os pedidos de clarificação também contam como recusa.

\subsection{Importância das Métricas de Recuperação}

Na maior parte dos sistemas RAG, os erros têm origem na pesquisa: se o excerto certo não é recuperado, nenhum modelo consegue responder corretamente. Por isso, a recuperação é avaliada antes da geração, com o script \texttt{check\_retrieval.py}, que mede, para cada pipeline e sem chamar nenhum LLM, o hit@k (alguma ata certa nos top-k) e o recall@k (fração das atas certas recuperadas, relevante nas perguntas multi\_ata), por tipo de pergunta.

Como a recuperação é idêntica para todos os modelos dentro de um pipeline, as métricas de contexto (context\_precision, context\_recall, retrieval\_hit@k) devem produzir resultados aproximadamente iguais entre modelos. Servem para:
\begin{itemize}
    \item Demonstrar que os erros vêm da geração e não da pesquisa
    \item Comparar as variantes de pesquisa (Secção~\ref{sec:pipelines})
    \item Validar que o pipeline de recuperação está a funcionar corretamente
\end{itemize}

\section{Conjunto de Perguntas}
\label{sec:perguntas}

\subsection{Categorias de Perguntas}

Foram definidos 8 tipos principais de perguntas, cada um testando capacidades específicas dos modelos:

\begin{table}[h!]
\centering
\caption{Categorias de perguntas e objetivos}
\label{tab:tipos-perguntas}
\begin{tabular}{>{\raggedright\arraybackslash}p{4cm}>{\raggedright\arraybackslash}p{5cm}>{\raggedright\arraybackslash}p{4cm}}
\toprule
\textbf{Tipo} & \textbf{Exemplo} & \textbf{O que testa} \\ \midrule
factual\_simples & Data da reunião da ata n.º 7 & Extração direta de factos \\ \midrule
votacao & Foi aprovada por unanimidade? Quem votou contra? & Interpretação de resultados de votação \\ \midrule
numerico & Valor do apoio à associação X & Precisão em números e valores monetários \\ \midrule
entidades & Com que entidades foi assinado o protocolo? & Identificação de nomes de pessoas/entidades \\ \midrule
temporal & Quando foi aprovado pela 1.ª vez X? & Raciocínio sobre datas e cronologia \\ \midrule
multi\_ata & Em quantas reuniões se discutiu X? & Agregação de informação de várias atas \\ \midrule
sintese & Resume as deliberações da reunião X & Capacidade de síntese e resumo fiel \\ \midrule
sem\_resposta & Orçamento do ``aeroporto de Barcelos'' & Resistência à alucinação \\ \bottomrule
\end{tabular}
\end{table}

\subsection{Perguntas Adicionais para Robustez}

Para além dos 8 tipos principais, foram adicionadas categorias para testar a robustez dos modelos face a:

\begin{itemize}
    \item \textbf{ruido\_reformulacao}: Perguntas com formulação alternativa mas mesmo conteúdo
    \item \textbf{resposta\_com\_justificacao}: Perguntas que exigem justificação da resposta
    \item \textbf{resistencia\_a\_premissa\_falsa}: Perguntas com premissas incorretas que o modelo deve corrigir
    \item \textbf{resistencia\_a\_entidade\_inexistente}: Perguntas sobre entidades que não existem nas atas
    \item \textbf{contradicao\_temporal}: Perguntas com informação temporal contraditória
    \item \textbf{resistencia\_a\_instrucao}: Perguntas que tentam forçar o modelo a ignorar as atas
    \item \textbf{resistencia\_a\_informacao\_injetada}: Perguntas com informação que o modelo deve verificar nas atas
    \item \textbf{distincao\_documental}: Perguntas que exigem distinção entre tipos de reunião (ordinária vs. extraordinária)
    \item \textbf{abstencao\_alucinacao}: Perguntas vagas onde o modelo deve recusar responder
    \item \textbf{resposta\_limitada\_a\_fonte}: Perguntas ambíguas onde o modelo deve pedir clarificação
\end{itemize}

\subsection{Implementação Atual}

Atualmente, o ficheiro questions.jsonl contém \textbf{28 perguntas} distribuídas por:
\begin{itemize}
    \item 3 perguntas factual\_simples (dificuldade: fácil)
    \item 2 perguntas votacao (dificuldade: média)
    \item 2 perguntas numerico (dificuldade: média/difícil)
    \item 2 perguntas entidades (dificuldade: média)
    \item 1 pergunta temporal (dificuldade: difícil)
    \item 3 perguntas multi\_ata (dificuldade: difícil)
    \item 1 pergunta sintese (dificuldade: média)
    \item 4 perguntas sem\_resposta (dificuldade: média/difícil)
    \item 10 perguntas de robustez, uma por categoria (vários níveis de dificuldade)
\end{itemize}

As respostas de referência foram escritas a partir do texto das atas e cada uma foi confirmada na ata de origem. Das 28 perguntas, 8 têm como resposta correta a recusa ou um pedido de clarificação.

\subsection{Estrutura dos Dados}

Cada pergunta no ficheiro questions.jsonl contém os seguintes campos:
\begin{itemize}
    \item \textbf{id}: Identificador único (ex: Q01, Q02, ...)
    \item \textbf{tipo}: Categoria da pergunta
    \item \textbf{dificuldade}: Nível de dificuldade (facil, media, dificil)
    \item \textbf{pergunta}: Texto da pergunta
    \item \textbf{resposta\_referencia}: Resposta correta escrita manualmente a partir das atas
    \item \textbf{valores\_esperados}: Lista de valores que devem estar presentes na resposta
    \item \textbf{ata\_fonte}: Lista de datas (AAAA-MM-DD) das atas que contêm a informação
\end{itemize}

\section{Implementação Técnica}
\label{sec:implementacao}

\subsection{Estrutura do Projeto}

\begin{verbatim}
atas-barcelos/
|-- config.yaml              # Pipelines, modelos e juiz
|-- .env.example             # Modelo do .env (chave pessoal)
|-- README.md                # Documentação do projeto
|-- requirements.txt         # Dependências Python
|-- data/
|   |-- atas/                # PDFs das atas
|   |-- questions.jsonl      # Perguntas de avaliação
|   `-- pipelines/<nome>/    # Artefactos por pipeline
|       |-- ocr/             # PDFs após OCR (cache)
|       |-- chunks.jsonl     # Chunks processados
|       `-- chroma/          # Índice vetorial ChromaDB
|-- src/
|   |-- common.py            # Config., cliente Ollama, embeddings
|   |-- ingest.py            # OCR, extração, metadados e chunking
|   |-- index.py             # Indexação com ChromaDB
|   |-- rag.py               # Retriever e Generator
|   |-- check_ollama.py      # Verifica a chave e o servidor
|   |-- check_retrieval.py   # Verifica a recuperação (sem LLM)
|   |-- run.py               # Execução dos modelos
|   `-- evaluate.py          # Avaliação dos resultados
|-- report/                  # Este relatório
`-- outputs/                 # Resultados
    |-- answers_<pipeline>__<modelo>.jsonl
    |-- resultados_por_pergunta.csv
    |-- resumo_por_modelo.csv
    `-- resumo_por_tipo.csv
\end{verbatim}

\subsection{Fluxo de Execução}

\begin{enumerate}
    \item \textbf{Ingestão} (ingest.py):
    \begin{itemize}
        \item OCR dos PDFs das atas, com cache dos PDFs reconhecidos
        \item Extração e limpeza de texto
        \item Identificação de metadados (número, data, tipo de reunião)
        \item Chunking por ``PROPOSTA N.º X''
        \item Gravação em data/pipelines/<pipeline>/chunks.jsonl
    \end{itemize}

    \item \textbf{Indexação} (index.py):
    \begin{itemize}
        \item Carregamento dos chunks
        \item Geração de embeddings com embeddinggemma, no servidor Ollama
        \item Criação do índice vetorial em ChromaDB, com o modelo de embeddings registado nos metadados
        \item Persistência em data/pipelines/<pipeline>/chroma/
    \end{itemize}

    \item \textbf{Verificação da recuperação} (check\_retrieval.py): hit@k e recall@k por tipo de pergunta, sem chamar nenhum LLM.

    \item \textbf{Execução} (run.py):
    \begin{itemize}
        \item Para cada pergunta, recuperação dos excertos (uma vez, partilhada por todos os modelos)
        \item Para cada modelo, geração da resposta com o contexto recuperado
        \item Gravação das respostas em answers\_<pipeline>\_\_<modelo>.jsonl
    \end{itemize}

    \item \textbf{Avaliação} (evaluate.py):
    \begin{itemize}
        \item Cálculo das métricas próprias (determinísticas)
        \item (Opcional) Cálculo das métricas RAGAS com LLM-juiz
        \item Geração de CSVs com resultados agregados por pipeline e modelo
    \end{itemize}
\end{enumerate}

\subsection{Detalhes de Implementação}

\paragraph{Ingestão (ingest.py)}
\begin{itemize}
    \item Aplica OCR com OCRmyPDF (Tesseract, português, correção de rotação)
    \item Utiliza PyMuPDF para extração de texto dos PDFs reconhecidos
    \item Trata particularidades das atas de Barcelos:
    \begin{itemize}
        \item Linhas de tracejado para preencher espaço
        \item Datas escritas por extenso
        \item Capitular na primeira letra
        \item Deliberações organizadas em ``PROPOSTA N.º X''
    \end{itemize}
    \item Converte datas por extenso para formato ISO (AAAA-MM-DD)
    \item Aplica limpeza: remove tracejados, números de página, normaliza espaços
\end{itemize}

\paragraph{Indexação (index.py)}
\begin{itemize}
    \item Utiliza ChromaDB como base de dados vetorial, embebida na aplicação
    \item Embeddings gerados no servidor Ollama da universidade (embeddinggemma, 768 dimensões), com o prefixo de documento recomendado para o modelo
    \item Índice persistente em disco, partilhado entre pipelines com o mesmo OCR, chunking e embeddings
    \item Medida de semelhança: cosseno
\end{itemize}

\paragraph{RAG (rag.py)}
\begin{itemize}
    \item \textbf{Retriever}: calcula o embedding da pergunta com o mesmo modelo usado na indexação (e recusa um índice feito com outro modelo) e devolve os top-k excertos; suporta pesquisa híbrida, filtro por data e reranking (Secção~\ref{subsec:pesquisa})
    \item \textbf{Generator}: envia a pergunta e os excertos ao modelo através do cliente Ollama autenticado, com repetição automática em caso de timeout
    \item \textbf{Prompt de sistema}: instruções para:
    \begin{itemize}
        \item Responder apenas com base nos excertos
        \item Usar ``A informação não consta das atas fornecidas.'' quando não há informação
        \item Responder em português europeu
        \item Indicar fontes no formato (Ata n.º X, AAAA-MM-DD)
    \end{itemize}
\end{itemize}

\paragraph{Avaliação (evaluate.py)}
\begin{itemize}
    \item Normalização de texto para comparação (remove acentos, normaliza números)
    \item Deteção automática de recusa e de pedidos de clarificação com base em padrões
    \item Cálculo de todas as métricas para cada pipeline x modelo x pergunta
    \item Agregação por pipeline, modelo e tipo de pergunta
    \item Exportação para CSV com pandas
\end{itemize}

\section{Executar o Pipeline}
\label{sec:execucao}

\subsection{Pré-requisitos}

O pipeline usa Python 3.11 e o Tesseract com a língua portuguesa:
\begin{verbatim}
sudo apt install ocrmypdf tesseract-ocr-por
conda create -n atas python=3.11 -y
conda activate atas
pip install -r requirements.txt
\end{verbatim}

Os modelos (geradores, juiz e embeddings) estão no servidor Ollama da universidade e não precisam de ser descarregados. Cada membro copia \texttt{.env.example} para \texttt{.env} e coloca a sua chave em \texttt{OLLAMA\_TOKEN}. O acesso verifica-se com \texttt{python src/check\_ollama.py}.

\subsection{Passos de Execução}

Os scripts aceitam \texttt{--pipeline <nome>}; sem esta opção, usam o \texttt{baseline}.

\begin{enumerate}
    \item Colocar os PDFs das atas em data/atas/
    \item Executar OCR e ingestão: python src/ingest.py
    \item Verificar o output para confirmar a extração correta de números e datas
    \item Executar indexação: python src/index.py
    \item Verificar a recuperação, sem LLM: python src/check\_retrieval.py
    \item Testar com limite: python src/run.py --limit 5
    \item Executar completo: python src/run.py
    \item Avaliar sem RAGAS (rápido): python src/evaluate.py --sem-ragas
    \item Avaliar com RAGAS (lento): python src/evaluate.py
\end{enumerate}

\section{Resultados Esperados}
\label{sec:resultados}

\subsection{Outputs Gerados}

O pipeline gera os seguintes ficheiros na pasta outputs/:
\begin{itemize}
    \item \textbf{answers\_<pipeline>\_\_<modelo>.jsonl}: Respostas brutas de cada modelo para cada pergunta
    \item \textbf{resultados\_por\_pergunta.csv}: Uma linha por pipeline x modelo x pergunta com todas as métricas
    \item \textbf{resumo\_por\_modelo.csv}: Médias de todas as métricas por pipeline e modelo
    \item \textbf{resumo\_por\_tipo.csv}: Médias por pipeline, modelo e tipo de pergunta
\end{itemize}

\subsection{Análise Planeada}

A análise dos resultados será focada em:
\begin{itemize}
    \item Tabela comparativa por modelo com todas as métricas
    \item Gráficos por tipo de pergunta para identificar pontos fortes e fracos
    \item Comparação das variantes do pipeline, em particular das variantes de pesquisa, primeiro com o check\_retrieval.py e depois com o melhor modelo
    \item Exemplos de erros típicos por modelo
    \item Correlação entre faithfulness (RAGAS) e valores\_corretos (determinística)
    \item Análise de custo-benefício: latência e tokens vs. qualidade das respostas
\end{itemize}

\section{Próximos Passos}
\label{sec:proximos}

\begin{enumerate}
    \item Concluir o OCR e a indexação do pipeline de referência
    \item Verificar a recuperação com o check\_retrieval.py e comparar as variantes de pesquisa
    \item Executar os seis modelos no pipeline de referência e avaliar os resultados
    \item Alargar o conjunto de perguntas (meta: 40--60 perguntas, $\geq$5 por tipo)
    \item Validar as respostas de referência com 2 pessoas do grupo
    \item Executar as restantes variantes do pipeline com o melhor modelo
\end{enumerate}

\section{Conclusões e Contribuições}
\label{sec:conclusoes}

Este projeto contribui com:
\begin{itemize}
    \item Um pipeline RAG padronizado para avaliação comparativa de LLMs
    \item Um conjunto de variantes do pipeline que permite medir o efeito de cada decisão de desenho
    \item Um conjunto de métricas abrangente (RAGAS + determinísticas)
    \item Um dataset de perguntas específicas para o domínio das atas municipais
    \item Decisões de desenho documentadas que garantem comparação justa
    \item Código reutilizável para futuras avaliações
\end{itemize}

As decisões de desenho, nomeadamente a recuperação fixa partilhada por todos os modelos, a variação de uma única decisão por variante do pipeline e a utilização de métricas determinísticas, permitem isolar os fatores de variação e obter conclusões robustas sobre o desempenho relativo dos diferentes modelos avaliados.

\end{document}
